%% kapa: exam=1 round=34 year=2023 period=2 form=A subject=appraisal_law subject_ko=감정평가관계법규 seq=23 number=23
%% source: Q-Net 감정평가사 기출문제 (cst00309) artlSeq=5211789 file=raw/1008/5211789/2023년도 제34회 감정평가사 1차시험 2교시.pdf
%% conversion: kordoc 4.24.51 -> tools/gen1.py (auto); review=proofread-sonnet; flags=none
\begin{question}{23}
국유재산법령상 일반재산에 관한 설명으로 옳지 않은 것은?
\begin{choices}
  \item 일반재산의 처분가격은 대통령령으로 정하는 바에 따라 시가(時價)를 고려하여 결정한다.
  \item 총괄청은 일반재산이 3년 이상 활용되지 않은 경우 이 일반재산을 민간사업자인 법인(외국법인 제외)과 공동으로 개발할 수 있다.
  \item 일반재산의 매각에 있어 매각 대금이 완납되기 전에 해당 매각재산의 소유권을 이전하는 경우에는 저당권 설정 등 채권의 확보를 위하여 필요한 조치를 취하여야 한다.
  \item 부동산신탁을 취급하는 신탁업자에게 신탁하여 개발된 일반재산의 대부기간은 30년 이내로 할 수 있으며, 20년의 범위에서 한 차례만 연장할 수 있다.
  \item 일반재산을 현물출자함에 있어 지분증권의 산정가액이 액면가에 미달하는 경우에는 그 지분증권의 액면가에 따라 출자가액을 산정한다.
\end{choices}
\end{question}
\answer{2}
